% Section 3: Phase 2 -- Q(sqrt(-5)) with class group + genus theory

Phase~2 implements the smallest non-trivial-class-number case
$K = \Q(\sqrt{-5})$, $h(K) = 2$. The Lemma~2.2 pigeonhole of
\cite{OpenAI2026Proof} becomes non-vacuous here: when an ideal product
$A_\varepsilon = \prod_j P_j^{\varepsilon_j} (P_j')^{k - \varepsilon_j}$
lands in the non-principal class, no element $\alpha_\varepsilon \in K^*$
with $(\alpha_\varepsilon) = A_\varepsilon A_\eta^{-1}$ exists for
arbitrary $\eta$ --- the pigeonhole must restrict to a single ideal
class, and the bound $|U| \geq (k+1)^s / h(K)$ is sharp.

\subsection{Classical facts encoded}

\begin{itemize}[leftmargin=*]
  \item $\OK = \Z[\sqrt{-5}]$ with discriminant $-20$
        (\texttt{K\_DISCRIMINANT = -20}, \texttt{CLASS\_NUMBER = 2}).
  \item The non-trivial ideal class is represented by
        $\mathfrak{p}_2 = (2, 1 + \sqrt{-5})$, satisfying
        $\mathfrak{p}_2^2 = (2)$.
  \item By genus theory \cite[Theorem~6.1]{Cox1989}, the Hilbert class
        field is the genus field $H = K(\sqrt{-1}) = \Q(i, \sqrt{5})$.
  \item A rational prime $p$ (odd, $p \neq 5$) splits completely in $K$
        iff $\left(\frac{-5}{p}\right) = 1$, equivalently
        $p \bmod 20 \in \{1, 3, 7, 9\}$. Of these, $p \bmod 20 \in \{1, 9\}$
        gives a principal prime (representable as $a^2 + 5b^2$), and
        $p \bmod 20 \in \{3, 7\}$ gives a non-principal prime.
\end{itemize}

\subsection{Lemma~2.2 pigeonhole, explicit}

The Phase~2 module \texttt{erdos\_ant.genus\_class\_field} exposes a
\texttt{GenusClassFieldProbe} that, given a tuple of split rational
primes, computes the ideal-class partition of the
$\{0,1\}^s$ family of ideal products and reports
$|U| \geq (k+1)^s / h(K)$.

\begin{verification}\label{ver:lemma22_principal}
For the choice $(p_1, p_2) = (29, 41)$ (both in the principal genus,
$29 = 3^2 + 5 \cdot 2^2$, $41 = 6^2 + 5 \cdot 1^2$), all four candidate
ideals $A_\varepsilon$ for $\varepsilon \in \{0,1\}^2$ are principal,
and the Lemma~2.2 lower bound is
$(1 + 1)^2 / 2 = 2$. Pinned in
\tname{tests/test_genus_class_field.py::test_lemma_22_pigeonhole_with_two_principal_primes}.
\end{verification}

\begin{verification}\label{ver:lemma22_mixed}
For the choice $(p_1, p_2) = (29, 3)$ (one principal, one
non-principal), the class group has order $2$, so the prime above $3$
and its conjugate lie in the same non-principal class. Consequently all
four ideals land in the non-principal class rather than splitting
$2 + 2$ between the two classes, and the Lemma~2.2 lower bound is still
$2$.
Pinned in \tname{test_lemma_22_pigeonhole_with_mixed_genus}.
\end{verification}

\subsection{Non-claims}

Phase~2 fixes a single CM field of degree $2$ over $\Q$ and produces a
finite pigeonhole demonstration. It does not produce a sequence of
fields with $[K_j : \Q] \to \infty$, and does not generate the
exponentially growing set $U$ that drives the asymptotic exponent
excess. It does, however, instantiate the smallest concrete example
where the pigeonhole step is genuinely non-trivial.
